\section{Methods}\label{sec:methods}

\subsection{Pipeline and stock-path generation}

The framework connected stock-path generation to contract-specific IV dynamics and American-option valuation. Information passed through the three stages in this order:
\begin{equation}\label{eq:pipeline}
    \text{JumpHMM stock path}
    \longrightarrow \text{contract IV factor}
    \longrightarrow \text{American lattice price and Greeks}.
\end{equation}
Alpaca option data entered during the offline surface fit. Each simulated step then used that fitted surface without requiring a new option quote or feeding lattice prices back into the fit. Stock paths came from pretrained single-ticker JumpHMM marginals fitted to equity histories from 2014--2024~\cite{jumphmm2025}. The model represents excess growth rates with latent states and location-scale Student-$t$ emissions. A jump mechanism increases visits to tail states to reproduce features such as heavy tails and volatility clustering~\cite{cont2001}. The pretrained artifact contained 424 ticker marginals with 50 states and 2{,}767 training days. Its growth rates used a $1/252$ year step and zero risk-free offset. A Student-$t$ copula can couple these marginals across assets~\cite{demarta2005}, although the reported scenarios used one ticker at a time.

For price simulation, we replaced the unbounded Student-$t$ emissions with their conditional distributions on a finite interval. For state $k$, the simulated excess growth rate was defined by:
\begin{equation}\label{eq:truncated_emission}
 G_t=\mu_k+\sigma_k Z_t,\qquad
 Z_t\sim t_{\nu_k}\mid |Z_t|\le b,\qquad b=10.
\end{equation}
The locations, scales, transition matrix, and jump parameters retained their training-fitted values. The cutoff was a simulation assumption applied after fitting; we did not refit the model under a truncated likelihood. A scale unit was the fitted Student-$t$ scale, not its standard deviation. We rejected and redrew out-of-bound emissions before applying the growth shift and converting log returns to prices. All stock paths were retained. This construction preserved the state and jump mechanism while giving finite stock-price moments at each fixed horizon. Without truncation, exponentiating Student-$t$ log returns gives an infinite stock-price mean and invalidates finite population interpretations of mean call values and short-call expected losses~\cite{cassidy2009}. We specified a wider cutoff $b=20$ to examine sensitivity, using the same contracts, seeds, and evaluation dates and recomputing normalization constants for each cutoff. Neither cutoff was selected by forecast performance. Supplementary Section~\ref{sec:truncated_emissions} reports the resulting bounds and sensitivity calculations.

\subsection{Offline IV-surface calibration}

We downloaded option-chain snapshots daily from Alpaca Markets using a free account. The account supplied Alpaca's indicative feed, whose quotes are modified derivatives of the exchange feed rather than actual OPRA quotes~\cite{alpaca_option_feed}. All option quote, midpoint, spread, and reported-IV comparisons in this study refer to these captured vendor values. They do not measure agreement with exchange-observed or executable prices. The frozen corpus contained 457 ticker files in fifteen snapshot directories and 234{,}549 observations after filtering. Calls and puts required a positive bid and positive calendar days to expiration (DTE). We also required reported IV strictly between 0.01 and 2.0 and strike/spot between 0.8 and 1.2 inclusive. Spot was the stored underlying-session close. Both option types entered the same loss, and repeated contracts across snapshots remained separate observations. The calibration target was Alpaca's reported IV field. Alpaca documents a Black--Scholes calculation with an iterative solver based on price sensitivity to volatility (Vega)~\cite{alpaca_market_data_faq}. Snapshot labels and underlying-session dates defined the respective temporal splits as detailed in Supplementary Table~\ref{tab:corpus_dates}. This corpus provided extensive cross-sectional coverage over a short observation window for comparing a five-parameter log-polynomial with neural surfaces shared globally, within sectors, or fitted separately to individual tickers. Parametric models such as SABR and SVI motivate compact representations of IV smiles~\cite{hagan2002,gatheral2004}. Neural representations allow more flexible shapes~\cite{horvath2021,bayer2019}. We exponentiated each shape function to keep its output positive. Neither calendar monotonicity nor strike convexity was imposed during fitting. The resulting fitted surfaces therefore required separate no-arbitrage checks. We used the following parametric shape for ticker $i$, calendar days to expiration $\tau$, and moneyness $m=K/S$:
\begin{equation}\label{eq:psi}
 \psi_{\beta}(\tau,m)=\exp\!\left[
 \beta_1\ln\tau+\beta_2\ln m+\beta_3\ln\tau\ln m
 +\beta_4(\ln m)^2+\beta_5(\ln\tau)^2\right].
\end{equation}
The corresponding IV prediction was defined by a ticker level and the surface shape:
\begin{equation}\label{eq:sigma_model}
    \widehat\sigma_i^2(\tau,m)=\theta_i\,\psi(\tau,m),
    \qquad \theta_i>0,\quad \psi(\tau,m)>0.
\end{equation}
The fitted squared IV also defined the variance target on each observation date:
\begin{equation}\label{eq:theta_calibration}
    \theta_{\mathrm{cal}}(i,\tau,m)=\theta_i\,\psi(\tau,m).
\end{equation}

The neural network (NN) models replaced the log-polynomial with a positive network output. We defined the two-input neural surface with the expression:
\begin{equation}\label{eq:psi_nn}
    \psi_{\mathrm{NN}}(\tau,m;w)
    =\exp\!\left[\mathrm{NN}\!\left(
    \operatorname{std}(\ln\tau),\operatorname{std}(\ln m);w\right)\right].
\end{equation}
Per-ticker fits required at least 2{,}000 observations; other tickers retained sector fits. Inputs were standardized using the training rows. Networks with at least 5{,}000 per-ticker observations, or 2{,}000 pooled group observations, used a $2\to16\to16\to1$ architecture with $\tanh$ activations and 337 parameters. Smaller qualified groups used $2\to8\to8\to1$ with 105 parameters. For the earnings analysis, we appended two standardized inputs: signed days to the ticker's nearest earnings date, clipped to $[-30,30]$, and the minimum absolute earnings distance among same-sector peers, clipped to 30. The resulting $4\to16\to16\to1$ network had 369 parameters. Earnings dates came from the public Yahoo Finance calendar. These event features entered only the temporal calibration comparison and remained inactive in the GS and LLY scenario drivers. To fit each network representation, we optimized its weights and ticker log levels jointly by minimizing the average squared difference between predicted and reported IV. Lower loss indicated closer agreement with the observed values. This criterion assigned equal weight to every retained observation. We minimized the following training objective:
\begin{equation}\label{eq:nn_loss}
 \min_{w,\{\ln\theta_i\}}
 \frac{1}{n}\sum_{j=1}^{n}
 \left[
 \sqrt{\theta_{i(j)}\psi_{\mathrm{NN}}(\tau_j,m_j;w)}
 -\sigma_{\mathrm{market},j}
 \right]^2.
\end{equation}
The log levels were initialized from each ticker's mean squared market IV. Adam training used a decreasing learning-rate schedule, at most 2{,}000 epochs, and a best-training-loss checkpoint with patience 200. Because a constant can be shifted between $\ln\theta_i$ and the output bias of the network, these two components are not separately identified. Only their product in~\eqref{eq:sigma_model} is interpretable without an additional normalization constraint.

\subsection{Contract-specific dynamic IV factors}\label{sec:theta_function}

The standard Heston model couples an asset-price diffusion to a square-root instantaneous-variance process~\cite{heston1993}. Its variance follows the stochastic differential equation:
\begin{equation}\label{eq:heston_standard}
    dv_t = \kappa(\theta-v_t)\,dt + \sigma_v\sqrt{v_t}\,dW_t^v .
\end{equation}
Here $\kappa$ is the mean-reversion rate, $\theta$ is long-run variance, and $\sigma_v$ is volatility of variance. The standard Heston model correlates the variance Brownian motion with the Brownian motion driving the stock. We used the square-root process to evolve a separate implied-variance state for each contract along a stock path supplied by JumpHMM. Each state reverted toward a target from the fitted surface and supplied the volatility input for that contract's lattice price. The lattice treated this input as constant within each pricing calculation even though it changed between simulated dates. This construction gave the square-root process the role of a dynamic IV factor. It did not represent the underlying asset's instantaneous return variance in a standard Heston economy. Within this construction, each scenario contract $a=(i,K,T)$ received a variance state $v_{a,t}$ with a target that changed as spot and remaining calendar DTE changed. Evaluating the fitted surface at the current contract state gave the target:
\begin{equation}\label{eq:theta}
    \bar v_{a,t}=\theta_i\,
    \psi_{\mathrm{NN}}\!\left(\tau_t,K/S_t\right).
\end{equation}
The factor started exactly on the fitted entry surface:
\begin{equation}\label{eq:v0}
    v_{a,0}=\bar v_{a,0}.
\end{equation}
At each later date, the lattice IV was the square root of that contract's factor:
\begin{equation}\label{eq:sigma_path}
    \sigma_{a,t}=\sqrt{v_{a,t}}.
\end{equation}

The put and call on a stock shared its path and the variance shock but retained their own variance states and surface targets. The stock and factor updates advanced only on exchange sessions. For the original fitted illustrations, we shifted every simulated growth-rate observation by a constant so that the ensemble mean matched a prespecified annual growth prior. We then standardized the one-session log returns across all paths and dates. The ablation used independent pilot estimates for these constants as described in Section~\ref{sec:ablation_methods}. We defined the standardized return shock entering the IV update with the expression:
\begin{equation}\label{eq:return_shock}
    Z^S_{j,t}=\frac{\Delta\ln S_{j,t}-\overline{\Delta\ln S}}
    {s_{\Delta\ln S}}.
\end{equation}
An independent standard normal shock supplied the remaining variation. Combining the return and independent shocks gave the factor shock:
\begin{equation}\label{eq:variance_shock}
    Z^v_{j,t}=\rho Z^S_{j,t}+\sqrt{1-\rho^2}\,Z^{\perp}_{j,t},
    \qquad Z^{\perp}_{j,t}\sim\mathcal N(0,1).
\end{equation}
The coefficient $\rho$ was fixed at $-0.6$. It was a return-shock coupling parameter, not a calibrated correlation between two Brownian motions in a standard Heston model.
We advanced each factor with a hard-floor Euler update and trading-year step $\Delta t=1/252$:
\begin{equation}\label{eq:euler}
 v_{a,t+1}=\max\!\left\{
 v_{a,t}+\kappa(\bar v_{a,t}-v_{a,t})\Delta t
 +\sigma_v\sqrt{\max(v_{a,t},0)}\sqrt{\Delta t}\,Z^v_{j,t},
 v_{\min}\right\}.
\end{equation}
We fixed $\kappa=15$, $\sigma_v=0.5$, and $v_{\min}=(0.005)^2$ as scenario assumptions. The floor protected the numerical pricer when the parameters did not satisfy the Feller condition~\cite{andersen2008}. The repository also contains a proposed state-dependent level and aggregate tail-state mood multiplier. Both terms remained inactive in every reported experiment. The fitted surface therefore determined the target throughout the study, and the fixed coefficient $\rho$ determined return-shock coupling. Evaluating the additional mechanisms would require state-conditioned option histories.

\subsection{American-option pricing, calendar clock, and terminal profit and loss}

Recombining lattices price American options by comparing continuation with immediate exercise at every node during backward induction~\cite{broadie1996}. The Cox--Ross--Rubinstein (CRR) construction sets its up and down factors from the supplied volatility~\cite{cox1979}. Alternatives include the moment matches of Jarrow--Rudd and Tian~\cite{jarrow1983,tian1993} and the Peizer--Pratt inversion used by Leisen--Reimer to improve finite-step convergence~\cite{leisen1996}. Both pricers used here accepted one IV value for a given contract and date. The GS and LLY contracts ran from April 28 to May 29, 2026. This interval contained 31 calendar days and 22 exchange-session transitions after excluding weekends and the May 25 market holiday. The surface and lattice used calendar DTE, with $\tau_t/365$ years entering the pricer. The stock and variance states advanced on the 22 trading transitions. Within this common clock, corpus-wide dollar-price checks used a 200-step CRR American tree. Scenario marks and Greeks used a 201-step Leisen--Reimer (LR) American tree. Both stages used a constant risk-free rate of $4.25\%$ and continuous dividend yield $q=0$. At expiration, the model replaced the lattice value with intrinsic value.

We calculated terminal profit and loss for a short position that received the captured indicative midpoint $p_0$ as follows:
\begin{equation}\label{eq:terminal_pnl}
 \Pi_T^{\mathrm{put}}=p_0^{\mathrm{put}}-\max(K-S_T,0),
 \qquad
 \Pi_T^{\mathrm{call}}=p_0^{\mathrm{call}}-\max(S_T-K,0).
\end{equation}
Equation~\eqref{eq:terminal_pnl} shows that terminal profit and loss depended on the JumpHMM terminal stock price, strike, and premium. Dynamic IV affected interim option marks and Greeks, but not a held-to-expiration payoff. All option premiums, values, and P\&L are dollars per underlying share; monetary amounts for a standard 100-share contract are 100 times these values. Premium financing was omitted. For low-volatility states or saturated Leisen--Reimer probabilities, an American CRR tree supplied the mark, retaining the immediate-exercise comparison.
We calculated price sensitivities, or Greeks, by central finite differences on the same Leisen--Reimer pricer. Delta measured price sensitivity to spot, and gamma measured the change in delta with spot. Spot bumps were $\pm1.5\%$ of spot for delta and gamma. Volatility bumps were $\pm0.005$, and Vega was reported per one volatility percentage point. We assumed execution at the observed mid, no transaction costs, no margin constraint, no early assignment before the plotted horizon, and no discrete dividends.

\subsection{Paired dynamic IV ablation}\label{sec:ablation_methods}

We tested which IV components changed interim valuation by holding the stock paths, contracts, entry premiums, and pricing assumptions fixed across five variants. Frozen IV retained each contract's fitted entry value of $v$. The direct surface set $v_{a,t}=\bar v_{a,t}$ at every date. Deterministic mean reversion used~\eqref{eq:euler} with $\sigma_v=0$. An uncoupled stochastic factor used $\sigma_v=0.5$ and $\rho=0$, and the full factor used $\sigma_v=0.5$ and $\rho=-0.6$. The last three variants retained $\kappa=15$. All variants started at the same fitted entry IV and reused the same independent normal innovation wherever a stochastic factor was active. The evaluation applied this paired design to three seeds per ticker, 20260429, 20260430, and 20260431, with 1{,}000 paths per seed. For this ablation, a separate 10{,}000-path pilot with seed 20260905 supplied a constant growth shift to the 10\% annual prior and the resulting one-session log-return mean and standard deviation. Those constants were frozen before generating evaluation paths. Unlike the original fitted illustrations, evaluation paths were therefore not recentered or standardized using their own future realizations. The pilot and evaluation simulations each used 22 transitions, and every return was applied once after the initial stock price.

The primary endpoint was liquidation after ten trading transitions, on May 12, with 17 calendar days remaining. Secondary endpoints used five, fifteen, and twenty transitions. The short-position liquidation P\&L and its paired difference from the direct-surface reference were given by the following expressions:
\begin{equation}\label{eq:ablation_pnl}
    \Pi_{a,t}^{(b,j)}=p_{a,0}-P_{a,t}^{(b,j)},
    \qquad D_{a,t}^{(b,j)}=\Pi_{a,t}^{(b,j)}-\Pi_{a,t}^{(\mathrm{surface},j)}.
\end{equation}
Here $b$ denotes the IV variant and $j$ the shared stock path. We reported mean P\&L, its standard deviation, the 5\% quantile, and 5\% expected shortfall, defined as mean P\&L at or below that quantile. The paired mean difference had Monte Carlo standard error $s_D/\sqrt{n}$, with smaller values indicating more precise simulation estimates of that mean. Mean absolute mark differences measured the size of pathwise changes even when positive and negative differences canceled in the mean. We reported each seed separately and pooled 3{,}000 paths per ticker. These uncertainty estimates condition on the fitted models and pilot constants. We also checked exact agreement of entry marks and terminal payoffs across variants and examined the numerical accuracy and strike consistency of the intervening prices. Numerical checks compared 201- and 401-step marks on the first twenty paths of every evaluation seed at transitions ten and twenty. A separate audit used eleven fixed strikes from 0.8 to 1.2 times entry spot, the first fifty paths of the first seed, and transitions zero, ten, and twenty. Calls and puts at the scenario expiry shared the same shocks. We checked intrinsic and upper price bounds, strike monotonicity, vertical-spread bounds, and convexity at both tree depths, counting violations larger than \$0.01 per share. Convexity violations were measured as the dollar excess above the neighboring linear interpolation. Entry was counted once per variant and option type. These are necessary strike conditions; the audit did not test calendar or dynamic no-arbitrage.

\subsection{Chronological surface and forecast evaluation}\label{sec:chronological_methods}

We extended the temporal evaluation in two stages to distinguish surface transfer from forecasting complete contract paths. The first stage used the completed expanding-window sector-network experiment on 58 capture dates through July 17, 2026. For each of 53 folds, the network fitted the first $k$ dates and predicted the next captured date, with $k=5,\ldots,57$. This experiment retained the original filters and capture-date labels, used two surface inputs, and calculated feature standardization from each training fold alone. Its 986{,}492 accepted observations contained no crossed or missing asks and no missing or nonpositive displayed bid or ask sizes under an independent reconstruction of the saved row counts. Since prediction used the later observation's stock price to calculate moneyness, the experiment measured conditional transfer of the surface across dates. It did not evaluate stock forecasts or the stochastic IV factor. The second stage fitted six sector networks once using 1{,}173{,}538 observations from 68 underlying sessions ending July 31. Their architectures and training schedule matched the existing two-input sector model, with seed 42 and a maximum of 2{,}000 epochs. Training-loss checkpoints, ticker levels, and training-only standardization were saved with the fitted models. The later sample contained 425{,}061 filtered IV observations across 23 sessions from August 4 through September 4. These observations supplied the final evaluation and did not select architectures, fitting windows, parameters, or seeds.

The second stage used an explicit capture audit before fitting or selecting contracts. The collector recorded retrieval times in UTC and stored the underlying-session close with each option row. We accepted retrievals after 20:00 UTC on that session and before 13:30 UTC on the next exchange session. This rule admitted overnight retrievals of a previous close while excluding twenty intraday repeat captures. We retained the latest eligible file per ticker and underlying session and one row per exact option symbol. Two-sided quotes required finite positive bids, noncrossed asks, and positive displayed sizes. Surface fitting additionally required the original IV and moneyness ranges and positive DTE, recomputed from expiration and the underlying session. Only files retrieved by the end of July 31 could enter training. Retrieval times were not exchange quote timestamps, so this audit did not establish quote freshness or exact stock--option synchronization. This limitation applied to the interpretation of both IV and dollar-price errors.

The forward test used GS and LLY at successive observed origins in August and early September, with planned horizons of one and five exchange sessions. Each origin supplied one put and one call with 25--45 calendar DTE. For each option type, we selected the expiration nearest 31 DTE and then the strike nearest $K/S=0.95$ for puts or $1.05$ for calls, breaking ties by earlier expiration and symbol. Selection used only the origin observations. We matched each selected symbol at its exact endpoint session and recorded missing endpoints without substituting another contract or capture date. The collection's rolling maturity buckets left none of these roughly monthly contracts observed five sessions later. We retained their one-session evaluation and added a separate shorter-maturity cohort with 10--21 DTE, choosing the expiration nearest 14 DTE with the same strike and tie rules. This availability-driven amendment preceded inspection of aggregate forecast errors and was recorded in the experiment protocol. The shorter cohort supplied the five-session case study; it did not replace the original contracts or validate their longer trajectories. All five IV variants began at the contract's observed origin IV. To compare their evolution without confounding the result with the fitted surface's initial level error, we rescaled the frozen variance target once at the origin:
\begin{equation}\label{eq:origin_anchor}
\bar v^{*}_{a,t}=\sigma_{a,0,\mathrm{obs}}^2
\frac{\bar v_a(S_t,\tau_t)}{\bar v_a(S_0,\tau_0)}.
\end{equation}
This anchoring defined a separate forecast experiment; the original sensitivity ablation retained its fitted initialization. Frozen IV held the common origin value, direct evaluation followed $\bar v^{*}_{a,t}$, and the three evolving factors used that anchored target with the previously specified parameters. Each origin used 1{,}000 JumpHMM paths and shared normal innovations across IV variants. The initial return state was drawn from the saved model's stationary distribution; recent returns did not update that state in the original forecast. An independent 10{,}000-path pilot per ticker supplied the growth shift to the existing 10\% annual prior and the one-session return normalization. Evaluation paths were not recentered. A zero-log-drift Gaussian random walk with the pilot return standard deviation provided a stock benchmark. Origin stock prices and IV updated as observations became available, but fitted parameters remained fixed. This was a retrospective evaluation at successive forecast origins, rather than a single forecast issued on July 31.

We scored forecast mean option values against the later bid--ask midpoints, using mean absolute error as the prespecified primary point metric. We retained that calculation and added median-based absolute errors to distinguish the choice of point summary from the IV comparison. Signed error, root mean squared error, and the fraction of predicted means inside the quoted spread supplied complementary price diagnostics. Distributional evaluation used the empirical continuous ranked probability score (CRPS)~\cite{gneiting2007}, which rewards concentration near the observed value while penalizing misplaced probability mass. For $N$ simulated values $P_j$ and observed midpoint $y$, we calculated the score as:
\begin{equation}\label{eq:forecast_crps}
\mathrm{CRPS}=\frac{1}{N}\sum_{j=1}^{N}|P_j-y|
-\frac{1}{2N^2}\sum_{j=1}^{N}\sum_{k=1}^{N}|P_j-P_k|.
\end{equation}
Smaller values indicate better distributional forecasts. We also reported central 90\% interval coverage and width. Scores were averaged across matched contracts within each ticker and origin before averaging dates, and stock forecast scores were reported separately. A diagnostic repeated IV evolution using the observed stock path only when every intervening session was available. It was explicitly conditional on realized stock information and therefore did not constitute a joint forecast. Repricing with the endpoint's reported IV in the same lattice supplied a further diagnostic of pricing discrepancies. All monetary scores were dollars per share. Overlapping horizons and common market movements made forecast dates and contracts dependent, so simulated path counts were not treated as independent market observations. The illustrative forecast was the earliest GS origin in the shorter-maturity cohort with both selected contracts observed at the five-session endpoint, regardless of its forecast errors.

\subsection{Stock forecast benchmarks}\label{sec:stock_validation_methods}

The option forecast errors motivated a separate evaluation of the stock generator. We first checked return scaling and price reconstruction and repeated the 2026 stock forecasts after conditioning the initial state on the available return history. This diagnostic retained the fitted parameters and tracked the remaining duration of a jump episode during state inference. We then compared four stock forecast methods: the saved 2014--2024 JumpHMM, an unchanged-price benchmark, a Gaussian model with adaptive volatility, and the same volatility model with a simple directional regression. The original JumpHMM retained its stationary initialization and independent-pilot growth shift. The adaptive model updated daily variance with an exponentially weighted moving average (EWMA) of squared log returns. The regression estimated a daily drift from the stock's most recent one, five, and twenty valid daily returns and SPY's most recent five valid daily returns. Cumulative-return features were normalized by origin volatility and the square root of their window length. A ridge penalty shrank the coefficients toward zero, and the regression had no intercept. Let $u_t=\log(S_t/S_{t-1})$ denote an observed daily return and $a_t$ its estimated daily variance at the forecast origin. The adaptive and directional models used the equations:
\begin{equation}\label{eq:stock_benchmark}
a_t=\lambda a_{t-1}+(1-\lambda)u_t^2,
\qquad
\log\frac{S_{t+h}}{S_t}=h\mu_t+\sqrt{a_t}\sum_{j=1}^{h}Z_j,
\quad Z_j\sim\mathcal N(0,1).
\end{equation}
The adaptive model set $\mu_t=0$, while the directional model used the regression estimate. Drift and variance remained fixed within each forecast horizon, so neither model received future stock or market returns. The unchanged-price benchmark supplied only a point forecast. It was not assigned an uncertainty interval for comparison with the stochastic models. These stock tests did not replace the paths used in the original option evaluation or the paired IV ablation.

We selected the new models' settings using expanding annual validation blocks from 2019 through 2024, with earlier observations used for each regression fit and feature scale. One EWMA decay and one ridge penalty were chosen across both tickers by weighting each ticker, validation year, and horizon equally. Candidate decays were $0.80$, $0.90$, $0.94$, $0.97$, and $0.99$. Ridge penalties ranged from $0.01$ to $1000$ and included the zero-drift limit. Decay selection minimized Gaussian log loss for cumulative log returns, while penalty selection minimized squared cumulative-return error divided by forecast variance. After selection, we fitted the regression using 2014--2024 observations and froze its coefficients and scaling. The additional 2025 stock dataset supplied 249 one-session and 245 five-session origins per ticker. The August--September 2026 test retained the original 21 one-session and 17 five-session stock outcomes per ticker. All methods used identical starting prices and endpoints. Origin features and variance updated only with observations then available; missing closes were not filled or converted into multi-session returns treated as one-day observations. Each stochastic model used 10{,}000 paths per origin and three fixed simulation seeds, with shared Gaussian innovations between the adaptive and directional candidates. We reported median-based MAE, mean-based RMSE, CRPS, and central 90\% interval coverage and width. Medians minimize expected absolute loss, while means minimize expected squared loss~\cite{hyndman2021}. Scores were averaged across dates within each simulation replicate and then across replicates. The model classes were proposed after inspecting the original 2026 misses, but neither the 2025 nor 2026 outcomes selected their settings. We therefore reported both periods separately and treated the 2026 follow-up as exploratory.

